% Figure: shear flow around an open channel section and the location of
% its shear center, offset from the web on the side opposite the
% flanges. A transverse load applied through the shear center bends
% without twisting; applied through the centroid it twists.
\begin{figure}[htbp]
\centering
\begin{tikzpicture}[scale=1.0]
% channel outline: web on the right, two flanges to the left
\draw[engblack, very thick] (3,2.2) -- (0,2.2);   % top flange
\draw[engblack, very thick] (3,2.2) -- (3,-2.2);  % web
\draw[engblack, very thick] (3,-2.2) -- (0,-2.2); % bottom flange
% shear flow arrows along the flanges (into the web)
\draw[engblue, very thick, ->] (0.3,2.2) -- (1.6,2.2);
\draw[engblue, very thick, ->] (1.6,2.2) -- (2.7,2.2);
\draw[engblue, very thick, ->] (0.3,-2.2) -- (1.6,-2.2);
\draw[engblue, very thick, ->] (1.6,-2.2) -- (2.7,-2.2);
% shear flow down the web
\draw[engblue, very thick, ->] (3,1.9) -- (3,0.2);
\draw[engblue, very thick, ->] (3,0.2) -- (3,-1.9);
\node[engblue, font=\scriptsize, anchor=south] at (1.3,2.25) {$q$};
\node[engblue, font=\scriptsize, anchor=west] at (3.05,1.0) {$q$};
% centroid C
\fill[enggray] (2.45,0) circle (0.07);
\node[enggray, font=\scriptsize, anchor=south] at (2.45,0.1) {$C$};
% shear center S, to the left of the web
\fill[engred] (0.9,0) circle (0.09);
\node[engred, font=\scriptsize, anchor=south] at (0.9,0.12) {$S$};
% offset dimension e
\draw[<->, engamber] (0.9,-2.6) -- (3,-2.6)
  node[midway, below, font=\scriptsize]{offset $e$};
% applied load through the shear center: bends without twisting
\draw[engred, very thick, ->] (0.9,2.9) -- (0.9,2.0);
\node[engred, anchor=south, font=\small] at (0.9,2.95) {$V$};
\end{tikzpicture}
\caption[Shear flow and the shear center of a channel
section]{Shear flow $q=VQ/I$ around an open channel under a transverse
shear $V$. The flow runs inward along both flanges and down the web,
and its resultant is a force in the web plus a couple from the flange
flows. Equating that couple to $V$ acting at an offset $e$ locates the
\emph{shear center} $S$, the point through which the load must pass for
the section to bend without twisting. For a channel $S$ lies outside
the web, on the side away from the flanges, a measurable distance $e$
from the web; the centroid $C$ sits inside the section. A load applied
through the centroid, as a designer might assume, produces a torque
$V\,e$ that twists the open section, whose torsional stiffness is small.
This is the reason channels and angles are braced against twist or
loaded through a point chosen to fall near the shear center.}
\label{fig:vol03ch08:shear-center}
\end{figure}
